\documentclass[10pt,a4paper]{amsart}
\usepackage[margin=19mm]{geometry}
\usepackage{amsmath,amssymb,mathtools}
\usepackage[scr=rsfs,cal=euler]{mathalfa}
\usepackage{xfrac}
\usepackage[unicode,hidelinks,bookmarks=false]{hyperref}
\newcommand{\ie}{i.e.\ }
\newcommand{\cf}{cf.\ }
\newcommand{\resp}{respectively\ }
\newcommand{\Subsection}{\subsection}
\providecommand{\nobreakdash}{\nobreak\hbox{-}}
\setlength{\emergencystretch}{2em}
\multlinegap0pt
\usepackage{enumerate}
\usepackage{amssymb}
\usepackage{mathtools}
\usepackage{xcolor}
\usepackage{tikz}
\usepackage{tcolorbox}
\usepackage{bm}
\usepackage{algorithm}\usepackage{algpseudocode}
\newtheorem{proposition}{Proposition}
\newcommand{\R}{\mathbb{R}}
\title{Zonoid volumes are not log-submodular}
\author{Ruben Skorupinski}
\date{Source: arXiv:2608.07702v1 (CC BY 4.0)}
\begin{document}
\maketitle

\noindent\emph{Source and licence.} Zonoid volumes are not log-submodular. Version arXiv:2608.07702v1 (CC BY 4.0), distributed by arXiv under the Creative Commons Attribution 4.0 International licence, http://creativecommons.org/licenses/by/4.0/, as recorded in the arXiv licence metadata for this exact version and shown on the abstract page https://arxiv.org/abs/2608.07702v1. Full licence notice: \texttt{licenses/arxiv-2608.07702-CC-BY-4.0.txt}.\par\smallskip

\noindent\textit{Editorial scope.} Counted unit: the article's proposition prop:unimodular\_zonos (source0.tex lines 262--264). The article's complete original proof of it is delivered with it (source0.tex lines 266--283, one \texttt{proof} environment of 18 lines). No mathematics is written, repaired or re-derived: the statement and the proof are the article's own lines, in the article's own notation, and no retained line is substituted or re-flowed.  No further source-local context is needed: every cross-reference of every retained line resolves inside the two delivered spans.  Nothing was omitted from any retained span, so the package carries no omission record.  No retained line contains a citation, so the wrapper carries no bibliography and no bibliography entry.  No retained line contains a figure environment, an image-inclusion command or drawing code, so no image is delivered and the package declares no asset dependency.  Editorial formatting only, in addition: an AMS \texttt{amsart} preamble that restates the article's own declarations and macros used by the retained lines, namely the environments \texttt{proposition} on separate counters, together with the standard packages those lines need.  The retained environments print in this document as Proposition 1 in order of appearance, whereas the article prints the same items with the numbering of its own full document.  The complete native source of the article as distributed in the arXiv e-print for this exact version is bundled unchanged as \texttt{sources/207102/source0.tex}.
\begin{proposition} \label{prop:unimodular_zonos}
    Let $Z$ be a zonotope generated by a rank-$d$ matrix $V \in \R^{d \times n}$ and let $w$, $w'$ be additional generators such that the matrix $(V|w|w')$ is unimodular. Then $|Z||Z+[0,w]+[0,w']| \leq |Z+[0,w]||Z+[0,w']|$. Moreover, equality holds if and only if $w\perp (VV^T)^{-1}w'$.
\end{proposition}

\begin{proof}
    Since all $(d\times d)$-sub-determinants $\det(V_B)$ of the unimodular matrix $V=(v_1,...,v_n)$ are in $\{-1,0,1\}$ the volume of the zonotope generated by $V$ is given by 
    \begin{align*}
        |Z| = \sum_{|B|=d} | \det(V_B)|  =  \sum_{|B|=d}\det(V_B)^2
    \end{align*} 
    By Cauchy-Binet we can thus write $|Z| = \det({\sum_i v_iv_i^T}) = \det(VV^T)$ and therefore, using the matrix-determinant lemma, we deduce that for an additional generator $w$ we get
    \begin{align*}
        |Z(V|w)|= \det\left( \sum_{i \in [n]}v_i v_i^T + ww^T \right) = \det\left( \sum_{i \in [n]}v_i v_i^T\right) \left(1 + w^T \left(\sum_{i \in [n]}v_i v_i^T\right)^{-1} w  \right).
    \end{align*}
    Note that the final expression is $ |Z(V)|(1+w^T(VV^T)^{-1}w)$ and hence the multiplicative increase in volume from adding the generator $w$ to $Z(V)$ is exactly the factor $1+w^T(VV^T)^{-1}w$. Similarly, when adding the same generator $w$ to the zonotope $Z(V|w')$ the relative increase is given by $1 + w^T(VV^T + w'w'^T)^{-1}w$. Finally, applying the Sherman–Morrison formula to $VV^T + w'w'^T$ and temporarily defining $M = VV^T$ for readability yields
    
    \begin{align*}
        &w^T(M + w'w'^T)^{-1}w \leq w^TM^{-1}w  \iff w^T\left(M^{-1}- \frac{M^{-1}w'w'^TM^{-1}}{1 + w'^TM^{-1}w'}\right)w \leq w^TM^{-1}w\\
        &\iff w^T\frac{M^{-1}w'w'^TM^{-1}}{1 + w'^TM^{-1}w'}w \geq 0 \iff (w'^TM^{-1}w)^2 \geq 0.
    \end{align*}

    Since the last inequality is true, with equality if and only if $w'^TM^{-1}w = 0$, the result follows.
\end{proof}

\end{document}
